\scanpage{105}
\mspage{50}
% archive sheet number: 104
$2^{'\circ}$) Dualement, la $W$-\uline{équivalence fine gauche} donne le $=$ résultat sur toutes les cat. du diagramme où elle est admise, et où elle ne coïncide pas avec la grossière, i.e. sur les 4 catégories

\[
 \mathrm{Cat}_{/S},\ \underline{\mathrm{Prop}}\,S,\ \underline{\mathrm{Cofib}}\,S,\ \underline{\mathrm{Cofibsc}}\,S.
\]

Il y a lieu de noter la catégorie de fractions commune par le sigle

\[
 \mathrm{HOT}_W(S^o)^o \quad(\text{ou aussi }\mathrm{HOT}^o_W(S))
\]

puisqu'elle est isomorphe à celle notée par ce sigle.

Si on veut des notations sans ambiguïté, désignant des cat. bien définies, et non pas seulement ``à équivalence près'', je choisirai

\[
 \begin{aligned}
 \mathrm{HOT}_W(S)&\overset{\mathrm{def}}{=}\underline{\mathrm{Fibsc}}(S)(W^d_S)_{\deleted{\unsure{-1}}}^{\unsure{+1|-1}}\\
 \text{d'où}\qquad \mathrm{HOT}_W(S^o)&\underset{\substack{\uparrow\\\text{isomorphisme}\\\text{de catégories}}}{\simeq}
 \bigl(\underline{\mathrm{Cofibsc}}\,(S)(W^g_S)_{\deleted{\unsure{g|-1}}}^{\unsure{+1|-1}}\bigr)^{\deleted{\unsure{o}}}
 \end{aligned}
\]

où $W^d_S,W^g_S$ désignent les ens. de flèches correspondant à la $W$-équivalence fine droite resp. gauche. Pour $\mathrm{HOT}^{\mathrm{lc}}_W(S)$, je préfère une définition aussi symétrique que possible en $S,S^o$, \deleted{je choisirai} \inserted{et de} \inserted{plus} une catégorie aussi petite que possible, ce qui me conduit à choisir

\[
 \mathrm{HOT}^{\mathrm{lc}}_W(S)\overset{\mathrm{def}}{=}(\underline{\mathrm{Parf}}\,S)W_S^{-1},
\]

où on a noté par $W_S$ les $W$-équivalences dans $\underline{\mathrm{Parf}}\,S$. (Inutile cette fois de préciser lesquelles -- toutes les quatre coïncident dans $\underline{\mathrm{Parf}}\,S$ !)

